\rhead{RB-EL4: Human-Robot Interaction (HRI)}
\phantomsection 
\section*{Human-Robot Interaction (HRI) (RB-EL4)}
\label{elec:RB_4}

\noindent \small{\hyperref[main:scheme]{$\leftarrow$ Back to Scheme}} \vspace{0.5cm}

\noindent \textbf{Credits:} 4 (3-1-0)

\subsection*{Theory}
\noindent\textbf{Unit 1} \\
\textit{HRI Fundamentals and Models:} HRI paradigms (physical, teleoperation, shared autonomy), Levels of autonomy (SAE 0-5 for robotics), Cognitive models (NASA TLX workload, Situation Awareness), Trust in automation (Lee \& See model), Uncanny valley effect, Robot appearance (anthropomorphic, zoomorphic, functional), Social robot design principles.

\vspace{0.3cm}

\noindent\textbf{Unit 2} \\
\textit{Natural User Interfaces:} Multimodal interaction (speech, gesture, gaze, haptics), Speech recognition (ASR: HMM, end-to-end DNNs), Natural Language Understanding (NLU: intent classification, slot filling), Dialogue management (POMDP, rule-based), Gesture recognition (MediaPipe, OpenPose), Gaze estimation and attention modeling.

\vspace{0.3cm}

\noindent\textbf{Unit 3} \\
\textit{Social and Affective HRI:} Emotion recognition (facial FER, physiological signals: EDA, HRV), Expressive behavior generation (gesture timing, facial animation), Rapport building strategies, Theory of Mind in robots, Personalization and user modeling, Cultural adaptation in HRI, Ethical considerations (deception, privacy, attachment).

\vspace{0.3cm}

\noindent\textbf{Unit 4} \\
\textit{Shared Autonomy and Assistive Robotics:} HRI in manipulation (shared control, intent prediction), Haptic feedback and teleoperation (passivity-based control), Assistive technologies (exoskeletons, wheelchairs), Legibility and predictability in motion, Human-aware navigation (social costs, pedestrian avoidance), Safety standards (ISO 13482, UL 1740).

\vspace{0.3cm}

\noindent\textbf{Unit 5} \\
\textit{HRI Evaluation and Applications:} Experimental design (within/between subjects, counterbalancing), Metrics (success rate, time-to-task, user satisfaction SUS/NASA-TLX), Field vs. lab studies, HRI in domains (healthcare, education, manufacturing, service robots), Emerging trends (AR/VR HRI, brain-computer interfaces), Accessibility and inclusive design.

\newpage

\subsection*{Practical/Simulations}
\noindent\textbf{Session 1: ROS2 Human Detection} \\
Focuses on YOLO + OpenCV for person tracking.

\vspace{0.2cm}

\noindent\textbf{Session 2: Speech Recognition Interface} \\
Focuses on Whisper ASR, intent classification with HuggingFace.

\vspace{0.2cm}

\noindent\textbf{Session 3: Gesture Recognition} \\
Focuses on MediaPipe hands/pose, custom classifier training.

\vspace{0.2cm}

\noindent\textbf{Session 4: Emotion Recognition} \\
Focuses on FER datasets, lightweight CNN deployment.

\vspace{0.2cm}

\noindent\textbf{Session 5: Dialogue System} \\
Focuses on Rasa NLU + dialogue policies.

\vspace{0.2cm}

\noindent\textbf{Session 6: Shared Autonomy Demo} \\
Focuses on admittance control, human force input.

\vspace{0.2cm}

\noindent\textbf{Session 7: Human-Aware Navigation} \\
Focuses on social navigation layers in ROS Nav2.

\vspace{0.2cm}

\noindent\textbf{Session 8: HRI User Study} \\
Focuses on SUS questionnaire, A/B testing design.

\vspace{0.2cm}

\noindent\textbf{Session 9: Multimodal Integration} \\
Focuses on fusion of speech/gesture/gaze inputs.

\vspace{0.2cm}

\noindent\textbf{Session 10: HRI Application Prototype} \\
Focuses on assistive robot with voice/gesture control.

\vspace{0.2cm}

\newpage
\rhead{Curriculum Schema}
